
%Commonsense reasoning is drawing increasing attention due to its importance in improving language understanding for machines.
\emph{\textbf{Warning}: this paper contains content that may
be offensive or upsetting.}

% Recently there has been a surge of interest in injecting information from commonsense knowledge bases (CSKB) into language models to improve their performance over downstream tasks. These CSKBs, however, are mostly human-generated and may reflect societal biases such as ``\emph{lawyers are dishonest}.'' It is important that these biases are not conflated with the notion of commonsense. We study this unaddressed yet important problem by first focusing on biases in two representative CSKBs, ConceptNet and GenericsKB.
% We identify two types of representational harms: \emph{overgeneralization} of polarized perceptions and representation \emph{disparity} across CSKBs and demographic groups.
% We find that these CSKBs contain severe overgeneralization and representation disparity issues.
% % across four demographic categories. 
% In addition, we analyze two downstream models that use ConceptNet as a source for commonsense knowledge and find similar representational harms from ConceptNet. We further propose a filtered-based approach and examine its effectiveness in reducing observed representational harms. We show that our filtered-based approach can reduce the issues in both resources and models but leads to a performance drop, leaving room for future work to build fairer and stronger commonsense models.
% \xiang{cut by 1/3}

Commonsense knowledge bases (CSKB) are increasingly used for various natural language processing tasks. Since CSKBs are mostly human-generated and may reflect societal biases, it is important to ensure that such biases are not conflated with the notion of commonsense. Here we focus on two widely used CSKBs, ConceptNet and GenericsKB, and establish the presence of bias in the form of two types of representational harms, \emph{overgeneralization} of polarized perceptions and representation \emph{disparity} across different demographic groups in both CSKBs. Next, we find similar representational harms for downstream models that use ConceptNet. Finally, we propose a filtering-based approach for mitigating such harms, and observe that our filtered-based approach can reduce the issues in both resources and models but leads to a performance drop, leaving room for future work to build fairer and stronger commonsense models.

%partial mitigation of bias in both resources and downstream models. 